\documentclass[12pt]{article}

\usepackage[a4paper, margin=2cm]{geometry}
\usepackage{setspace}
\usepackage{tabularx, booktabs}
\usepackage{caption}
\usepackage{graphicx}
\usepackage{amsmath}
\usepackage{minted}
\usepackage{float}
\usepackage{animate}

\usepackage{polyglossia}
\setdefaultlanguage{russian}
\setotherlanguages{english}

\usepackage{fontspec}
\setmainfont{NewComputerModern10}
\setsansfont{NewComputerModernSans10}
\setmonofont[
  ItalicFont={Monaspace Radon}
]{Monaspace Neon}

\newcolumntype{Y}{>{\centering\arraybackslash}X}

\begin{document}
\begin{titlepage}
\centering

{Министерство науки и высшего образования Российской Федерации}
\vskip 2mm
{\large\scshape Новосибирский государственный технический университет}
\vskip 2mm
{Кафедра прикладной математики}

\vfill
\vskip 3cm

{\scshape Курсовой проект}
\vskip 3mm
\begin{spacing}{1.8}
    {\huge\bf Метод конечных элементов для параболических уравнений}
\end{spacing}
\vskip 5mm
{\scriptsize ФПМИ, ПМ-24}
\vskip 2mm
{\large\it Параскун И.}

\vfill

{преподаватели}
\vskip 2mm
{\scshape Вагин Д. В., д.т.н., доцент}
\vskip 1mm
{\scshape Сивенкова А. П., ассистент}

\vfill

{Новосибирск, \the\year}

\end{titlepage}

\setcounter{page}{2}
\tableofcontents
\newpage

\section{Постановка задачи}

\noindent Построить МКЭ для уравнения параболического типа

\begin{align}
  \label{eq:1} -\nabla(\lambda \nabla u) + \gamma u + 
    \sigma\frac{\partial u}{\partial t} = f
\end{align}

\vspace{4mm}
\noindent в декартовой системе координат

\begin{align}
  - \frac{\partial}{\partial x}(\lambda\frac{\partial u}{\partial x}) 
  - \frac{\partial}{\partial y}(\lambda\frac{\partial u}{\partial y}) 
  - \frac{\partial}{\partial z}(\lambda\frac{\partial u}{\partial z})
    + \gamma u + \sigma\frac{\partial u}{\partial t} = f
\end{align}

\vskip 4mm
\noindent с учётом следующих краевых условий:

\begin{align}
  &u \big|_{S_1} = u_g \\
  &\lambda\frac{\partial u}{\partial n} \bigg|_{S_2} = \theta \\
  &\lambda\frac{\partial u}{\partial n} \bigg|_{S_3} +
    \beta(u \big|_{S_3} - u_{\beta}) = 0
\end{align}

\section{Теоретическая часть}

\subsection{Вариационная постановка}

\noindent Эквивалентная вариационная постановка в форме уравнения Галёркина для
эллиптического уравнения, входящего в уравнение \ref{eq:1}:

\begin{align} \label{eq:2}
  \begin{split}
    \int_{\Omega}\lambda \nabla u \nabla v_0 d\Omega &+
    \int_{\Omega}\gamma uv_0 d\Omega +
    \int_{S_3}\beta uv_0 dS \\
    &= \int_{\Omega}f v_0 d\Omega +
    \int_{S_2}\theta v_0 dS +
    \int_{S_3}\beta u_{\beta} v_0 dS, \ \ \forall v_0 \in H_0^1
  \end{split}
\end{align}

\vskip 5mm
\noindent Аппроксимация уравнения \ref{eq:2} на конечномерных подпространствах $V^h_g$ и $V^h_0$ 
получается заменой функций $u \in H^1_g$ и $v_0 \in H^1_0$ на 
функции $u_h \in V^h_g$ и $v^h_0 \in V^h_0$ соответственно:

\begin{align}
  \begin{split}
    \int_{\Omega}\lambda \nabla u^h \nabla v^h_0 d\Omega &+
    \int_{\Omega}\gamma u^h v^h_0 d\Omega +
    \int_{S_3}\beta u^h v^h_0 dS \\
    &= \int_{\Omega}f v^h_0 d\Omega +
    \int_{S_2}\theta v^h_0 dS +
    \int_{S_3}\beta u_{\beta} v^h_0 dS, \ \ \forall v^h_0 \in V^h_0
  \end{split}
\end{align}

\vskip 5mm

\subsection{Конечноэлементная СЛАУ}

\noindent Раскладывая функции $u$ и $v_0$ по базису, переходим к конечноэлементной СЛАУ

\begin{align}
  \begin{split}
    \sum_{j=1}^{n}\big(
      \int_{\Omega}\lambda \nabla \psi_j \nabla \psi_i d\Omega &+
      \int_{\Omega}\gamma \psi_j \psi_i d\Omega +
      \int_{S_3}\beta \psi_j \psi_i dS
    \big)q_j \\
    &= \int_{\Omega}f \psi_i d\Omega +
    \int_{S_2}\theta \psi_i dS +
    \int_{S_3}\beta u_{\beta} \psi_i dS, \ \ i \in N_0
  \end{split}
\end{align}

\vskip 5mm
\noindent или в матричном виде

\begin{align}
  (\textbf{G} + \textbf{M}^{\gamma} + \textbf{M}^{S_3})\textbf{q} = \textbf{b}
\end{align}

\vskip 5mm
\noindent где компоненты определяются соотношениями

\begin{align}
  &G_{ij} = \int_{\Omega}\lambda \nabla \psi_j \nabla \psi_i d\Omega
          = \sum_{k}\int_{\hat{\Omega}}\hat{\lambda}\nabla\hat{\psi}_j\nabla\hat{\psi}_i d\hat{\Omega} \\
  &M^{\gamma}_{ij} = \int_{\Omega}\gamma \psi_j \psi_i d\Omega
                   = \sum_{k}\int_{\hat{\Omega}}\hat{\gamma}\hat{\psi}_j\hat{\psi}_i d\hat{\Omega}\\
  &M^{S_3}_{ij} = \int_{S_3}\beta \psi_j \psi_i dS
                = \sum_{k}\int_{\hat{S}_3}\hat{\beta}\hat{\psi}_j\hat{\psi}_i d\hat{S} \\
  &b^{\Omega}_{i} = \int_{\Omega}f \psi_i d\Omega
                  = \sum_{k}\int_{\hat{\Omega}}\hat{f}\hat{\psi}_i d\hat{\Omega} \\
  &b^{S_2}_{i} = \int_{S_2}\theta \psi_i dS
               = \sum_{k}\int_{\hat{S}_2}\hat{\theta}\hat{\psi}_i d\hat{S} \\
  &b^{S_3}_{i} = \int_{S_3}\beta u_{\beta} \psi_i dS
               = \sum_{k}\int_{\hat{S}_3}\hat{\beta}\hat{u}_{\beta}\hat{\psi}_i d\hat{S}
\end{align}

\vskip 5mm

\subsection{Трилинейные базисные функции}

\vskip 5mm
\noindent Трилинейные бизисные функции на прямоугольном параллелепипеде $\Omega_{psr}$ строятся
  следующим образом:

\begin{align*}
  & X_1(x) = \frac{x_{p+1}-x}{h_x}, & &X_2(x) = \frac{x - x_{p}}{h_x}, & &h_x = x_{p+1} - x_p, \\
  & Y_1(y) = \frac{y_{s+1}-y}{h_y}, & &Y_2(y) = \frac{y - y_{s}}{h_y}, & &h_y = y_{s+1} - y_s, \\
  & Z_1(z) = \frac{z_{r+1}-z}{h_z}, & &Z_2(z) = \frac{z - z_{r}}{h_z}, & &h_z = z_{r+1} - z_r.
\end{align*}

\begin{align}
  \label{eq:3} \hat{\psi}_i = X_{\mu(i)}Y_{\nu(i)}Z_{\vartheta(i)}
\end{align}

\vskip 5mm
\noindent Выражения для вычисления компонент локальных матриц и векторов на прямоугольном
параллелепипеде $\Omega_{psr}$:

\begin{align}
  \begin{split}
    \hat{G}_{ij} = \hat{\lambda}(
      &G^{x}_{\mu(k)\mu(j)\mu(i)}M^{y}_{\nu(k)\nu(j)\nu(i)}M^{z}_{\vartheta(k)\vartheta(j)\vartheta(i)} + \\
      &M^{x}_{\mu(k)\mu(j)\mu(i)}G^{y}_{\nu(k)\nu(j)\nu(i)}M^{z}_{\vartheta(k)\vartheta(j)\vartheta(i)} + \\
      &M^{x}_{\mu(k)\mu(j)\mu(i)}M^{y}_{\nu(k)\nu(j)\nu(i)}G^{z}_{\vartheta(k)\vartheta(j)\vartheta(i)}
    )
  \end{split}
\end{align}

\begin{align}
  &\hat{M}^{\gamma}_{ij} = \hat{\gamma}(
      M^{x}_{\mu(j)\mu(i)}M^{y}_{\nu(j)\nu(i)}M^{z}_{\vartheta(j)\vartheta(i)}
  ) \\
  &\hat{M}^{S_3}_{ij} = \hat{\beta}(
      M^{\xi}_{\mu(j)\mu(i)}M^{\chi}_{\nu(j)\nu(i)}
  ) \\
  &\hat{b}^{\Omega}_{i} = \sum_{k=1}^{8}\hat{f}_{k}(
      M^{x}_{\mu(k)\mu(i)}M^{y}_{\nu(k)\nu(i)}M^{z}_{\vartheta(k)\vartheta(i)}
  ) \\
  &\hat{b}^{S_2}_{i} = \sum_{k=1}^{4}\hat{\theta}_{k}(
      M^{\xi}_{\mu(k)\mu(i)}M^{\chi}_{\nu(k)\nu(i)}
  ) \\
  &\hat{b}^{S_3}_{i} = \hat{\beta}\sum_{k=1}^{4}\hat{u}_{\beta, k}(
      M^{\xi}_{\mu(k)\mu(i)}M^{\chi}_{\nu(k)\nu(i)}
  )
\end{align}

\vskip 5mm
\noindent где

\begin{align}
  &\textbf{G}^{\xi} = \frac{1}{h_{\xi}}
    \begin{pmatrix}
      1 & -1 \\
      -1 & 1
    \end{pmatrix} \\
  &\textbf{M}^{\xi} = \frac{h_{\xi}}{6}
    \begin{pmatrix}
      2 & 1 \\
      1 & 2
    \end{pmatrix}
\end{align}

\vskip 5mm

\subsection{Аппроксимация по времени}

\subsubsection{Двухслойная схема}

\noindent Полагая, что коэффициенты $\lambda$ и $\gamma$ не зависят от времени, неявная 
  двухслойная схема аппроксимации уравнения \ref{eq:1} по времени может быть выписана как

\begin{align}
  -\nabla(\lambda \nabla u^j) + \gamma u^j + 
    \sigma\frac{u^j - u^{j-1}}{\Delta t} = f^j, \ \ j = 1...J
\end{align}

\vspace{5mm}
\noindent где $u^{j-1}$ - решение, полученное на предыдущем временном слое, а $u^0$ - начальное условие.

\vspace{2mm}
\noindent Проводя преобразования, аналогичные тем, что были проведены для эллиптической составляющей, 
  получаем следующее матричное уравнение:

\begin{align}
  (\textbf{G} + \textbf{M}^{\gamma} + \textbf{M}^{S_3} + \frac{1}{\Delta t}\textbf{M}^{\sigma})\textbf{q}^j =
    \textbf{b}^j + \frac{1}{\Delta t}\textbf{M}^{\sigma}\textbf{q}^{j-1}
\end{align}

\vspace{3mm}

\subsubsection{Трехслойная схема}

\noindent Представим искомое решение $u$ на интервале $(t_{j-2}, t_j)$ в следующем виде:

\begin{align}
  \label{eq:4} u(x, y, z, t) \approx u^{j-2}(x, y, z)\eta^j_2(t) + 
                        u^{j-1}(x, y, z)\eta^j_1(t) +
                        u^{j}(x, y, z)\eta^j_0(t)
\end{align}

\vspace{5mm}
\noindent где функции $\eta^j_2(t), \eta^j_1(t), \eta^j_0(t)$ - базисные квадратичные полиномы
  Лагранжа, которые могут быть записаны в виде

\begin{align}
  &\eta^j_2(t) = \frac{(t - t_{j-1})(t - t_j)}{2\Delta t^2} \\
  &\eta^j_1(t) = -\frac{(t - t_{j-2})(t - t_j)}{\Delta t^2} \\
  &\eta^j_0(t) = \frac{(t - t_{j-2})(t - t_{j-1})}{2\Delta t^2}
\end{align}

\vspace{5mm}
\noindent Производные по времени на $j$-ом временном слое:

\begin{align}
  &\frac{d\eta^j_2(t)}{dt}\bigg|_{t = t_j} = \frac{1}{2 \Delta t} \\
  &\frac{d\eta^j_1(t)}{dt}\bigg|_{t = t_j} = -\frac{2}{\Delta t} \\
  &\frac{d\eta^j_0(t)}{dt}\bigg|_{t = t_j} = \frac{3}{2 \Delta t}
\end{align}

\vspace{5mm}
\noindent Применяя выражение \ref{eq:4} для аппроксимации производной по времени уравнения \ref{eq:1}
  на $j$-ом временном слое:

\begin{align}
  \begin{split}
    \frac{\partial}{\partial t}(
      &u^{j-2}(x, y, z)\eta^j_2(t) + 
      u^{j-1}(x, y, z)\eta^j_1(t) +
      u^{j}(x, y, z)\eta^j_0(t)
    )\bigg|_{t = t_j} \\ 
    &- \nabla(\lambda \nabla u^j) + \gamma u^j = f^j, \ \ j = 2...J
  \end{split}
\end{align}

\vspace{5mm}
\noindent или, с учётом выражений для вычисления производных

\begin{align}
  -\nabla(\lambda \nabla u^j) + \gamma u^j + 
    \frac{1}{2 \Delta t} u^{j-2} -
    \frac{2}{\Delta t} u^{j-1} + 
    \frac{3}{2 \Delta t} u^j = f^j
\end{align}

\vspace{5mm}
\noindent Проводя преобразования, аналогичные тем, что были проведены для эллиптической составляющей, 
  получаем следующее матричное уравнение:

\begin{align}
  (\textbf{G} + \textbf{M}^{\gamma} + \textbf{M}^{S_3} + 
    \frac{3}{2 \Delta t}\textbf{M}^{\sigma})\textbf{q}^j = \textbf{b}^j 
      - \frac{1}{2 \Delta t}\textbf{M}^{\sigma}\textbf{q}^{j-2}
      + \frac{2}{\Delta t}\textbf{M}^{\sigma}\textbf{q}^{j-1}
\end{align}

\vspace{3mm}

\subsubsection{Четырехслойная схема}

\noindent Представим искомое решение $u$ на интервале $(t_{j-2}, t_j)$ в следующем виде:

\begin{align}
  \label{eq:5} u(x, y, z, t) \approx u^{j-3}(x, y, z)\eta^j_3(t) +
                        u^{j-2}(x, y, z)\eta^j_2(t) + 
                        u^{j-1}(x, y, z)\eta^j_1(t) +
                        u^{j}(x, y, z)\eta^j_0(t)
\end{align}

\vspace{5mm}
\noindent где функции $\eta^j_3(t), \eta^j_2(t), \eta^j_1(t), \eta^j_0(t)$ - базисные кубические полиномы
  Лагранжа, которые могут быть записаны в виде

\begin{align}
  &\eta^j_3(t) = -\frac{(t - t_{j-2})(t - t_{j-1})(t - t_{j})}{6\Delta t^3} \\
  &\eta^j_2(t) = \frac{(t - t_{j-3})(t - t_{j-1})(t - t_{j})}{2\Delta t^3} \\
  &\eta^j_1(t) = -\frac{(t - t_{j-3})(t - t_{j-2})(t - t_{j})}{2\Delta t^3} \\
  &\eta^j_0(t) = \frac{(t - t_{j-3})(t - t_{j-2})(t - t_{j-1})}{6\Delta t^3}
\end{align}

\vspace{5mm}
\noindent Производные по времени на $j$-ом временном слое:

\begin{align}
  &\frac{d\eta^j_3(t)}{dt}\bigg|_{t = t_j} = -\frac{1}{3 \Delta t} \\
  &\frac{d\eta^j_2(t)}{dt}\bigg|_{t = t_j} = \frac{3}{2 \Delta t} \\
  &\frac{d\eta^j_1(t)}{dt}\bigg|_{t = t_j} = -\frac{3}{\Delta t} \\
  &\frac{d\eta^j_0(t)}{dt}\bigg|_{t = t_j} = \frac{11}{6 \Delta t}
\end{align}

\vspace{5mm}
\noindent Применяя выражение \ref{eq:5} для аппроксимации производной по времени уравнения \ref{eq:1}
  на $j$-ом временном слое:

\begin{align}
  \begin{split}
    \frac{\partial}{\partial t}(
      &u^{j-3}(x, y, z)\eta^j_3(t) + 
      u^{j-2}(x, y, z)\eta^j_2(t) + 
      u^{j-1}(x, y, z)\eta^j_1(t) +
      u^{j}(x, y, z)\eta^j_0(t)
    )\bigg|_{t = t_j} \\ 
    &- \nabla(\lambda \nabla u^j) + \gamma u^j = f^j, \ \ j = 2...J
  \end{split}
\end{align}

\vspace{5mm}
\noindent или, с учётом выражений для вычисления производных

\begin{align}
  -\nabla(\lambda \nabla u^j) + \gamma u^j - 
    \frac{1}{3 \Delta t} u^{j-3} +
    \frac{3}{2 \Delta t} u^{j-2} -
    \frac{3}{\Delta t} u^{j-1} + 
    \frac{11}{6 \Delta t} u^j = f^j
\end{align}

\vspace{5mm}
\noindent Проводя преобразования, аналогичные тем, что были проведены для эллиптической составляющей, 
  получаем следующее матричное уравнение:

\begin{align}
  (\textbf{G} + \textbf{M}^{\gamma} + \textbf{M}^{S_3} + 
    \frac{11}{6 \Delta t}\textbf{M}^{\sigma})\textbf{q}^j = \textbf{b}^j 
      + \frac{1}{3 \Delta t} \textbf{M}^{\sigma}\textbf{q}^{j-3}
      - \frac{3}{2 \Delta t} \textbf{M}^{\sigma}\textbf{q}^{j-2}
      + \frac{3}{\Delta t} \textbf{M}^{\sigma}\textbf{q}^{j-1}
\end{align}

\vspace{3mm}

\section{Описание разработанных программ}

\selectlanguage{english}
\inputminted[firstline=7, lastline=53]{c}{pde/include/numx/pde/msh.h}
\inputminted[firstline=11, lastline=80]{c}{pde/include/numx/pde/cnd.h}
\inputminted[firstline=8, lastline=139]{c}{pde/include/numx/pde/sim.h}
\inputminted[firstline=34, lastline=164]{c}{pde/include/numx/pde/slv.h}
\inputminted[firstline=6, lastline=29]{c}{pde/include/numx/pde/fem.h}
\selectlanguage{russian}

\newpage

\section{Описание тестирования программы}

\subsection{Тестирование на регурярной сетке}
\subsubsection{Описание задачи}

\begin{figure}[H]
  \centering
  \vspace{5mm}
  \includegraphics[scale=0.8]{pde/test/pbc/c1/img/geo*.png}
  \vspace{5mm}
  \caption{Засчётная область}
\end{figure}

\paragraph{Условия}
\indent
\vspace{1mm}

\noindent Решение - полином вида

\begin{align}
  u = 5xt^3 - 2zt^2 + yt + 1
\end{align}

\vspace{5mm}
\noindent Принимая значения коэффициентов 
  $\lambda = 1$, $\gamma = 0$, $\sigma = 2$, $\beta = 10$, получаем

\begin{align*}
  & f = 30xt^2 - 8zt + 2y \\
  & \theta|_{S_1} = -5t^3 \\
  & \theta|_{S_3} = t \\
  & \theta|_{S_5} = -2t^2 \\
  & u_{\beta}|_{S_2} = 5xt^3 - 2zt^2 + yt + 1 - 0.1t \\
  & u_{\beta}|_{S_4} = 5xt^3 - 2zt^2 + yt + 1 + 0.5t^3 \\
  & u_{\beta}|_{S_6} = 5xt^3 - 2zt^2 + yt + 1 + 0.2t^2
\end{align*}

\vspace{5mm}
\noindent Значения параметров по времени: $t_0 = 0$, $t_1 = 25$, $\Delta t = 1$.

\subsubsection{Результат выполнения}

\noindent Для деменстрации работы многослойных схем аппроксимации по времени,
  решения на первых временных слоях были взяты аналитически (1-3 начальных условия).

\begin{table}[H]
  \begin{tabularx}{\textwidth}{cYYY}
    \toprule
      $j$ & \multicolumn{3}{c}{$E(|u - u^*|)$}                              \\
    \cmidrule(lr){2-4}
          & Двухслойная схема   & Трёхслойная схема  & Четырёхслойная схема \\
    \midrule
      0	  & 0.0000000E+00	      & 0.0000000E+00	     & 0.0000000E+00        \\
      1	  & 1.3008431E+02	      & 0.0000000E+00	     & 0.0000000E+00        \\
      2	  & 3.9599125E+02	      & 1.0564511E+02	     & 0.0000000E+00        \\
      3	  & 7.2839704E+02	      & 1.9439680E+02	     & 3.7414677E-11        \\
      4	  & 1.0920167E+03	      & 2.4007804E+02	     & 1.8934309E-11        \\
      5	  & 1.4702653E+03	      & 2.5755264E+02	     & 5.3405582E-11        \\
      6	  & 1.8553821E+03	      & 2.6224142E+02	     & 7.1549513E-11        \\
      7	  & 2.2437265E+03	      & 2.6252408E+02	     & 1.9607504E-10        \\
      8	  & 2.6335882E+03	      & 2.6184524E+02	     & 1.2670078E-10        \\
      9	  & 3.0241632E+03	      & 2.6127019E+02	     & 3.8652964E-10        \\
      10	& 3.4150735E+03	      & 2.6096067E+02	     & 5.1888799E-10        \\
      11	& 3.8061415E+03	      & 2.6083480E+02	     & 6.3853943E-10        \\
      12	& 4.1972836E+03	      & 2.6079814E+02	     & 1.2761810E-09        \\
      13	& 4.5884605E+03	      & 2.6079426E+02	     & 1.0528181E-09        \\
      14	& 4.9796539E+03	      & 2.6079828E+02	     & 1.3561256E-09        \\
      15	& 5.3708549E+03	      & 2.6080206E+02	     & 1.8517922E-09        \\
      16	& 5.7620596E+03	      & 2.6080417E+02	     & 1.8093863E-09        \\
      17	& 6.1532660E+03	      & 2.6080504E+02	     & 2.1203369E-09        \\
      18	& 6.5444732E+03	      & 2.6080531E+02	     & 2.9247722E-09        \\
      19	& 6.9356808E+03	      & 2.6080535E+02	     & 2.6484560E-09        \\
      20	& 7.3268886E+03	      & 2.6080533E+02	     & 3.3724367E-09        \\
      21	& 7.7180964E+03	      & 2.6080530E+02	     & 3.8412062E-09        \\
      22	& 8.1093043E+03	      & 2.6080529E+02	     & 8.4636231E-09        \\
      23	& 8.5005122E+03	      & 2.6080528E+02	     & 5.3736107E-09        \\
      24	& 8.8917201E+03	      & 2.6080528E+02	     & 1.0244272E-08        \\
      25	& 9.2829280E+03	      & 2.6080528E+02	     & 8.3934569E-09        \\
    \bottomrule
  \end{tabularx}

  \captionsetup{justification=centering}
  \caption{Погрешность решения на $t_j$ временном слое}
\end{table}

\begin{figure}[H]
  \centering
  \animategraphics[autoplay,loop,width=0.8\textwidth]{15}{pde/test/pbc/c1/img/vis/vis.}{0000}{0025}
  \vspace{2mm}
  \caption{Визуализация решения}
\end{figure}

\subsection{Тестирование на порядок аппроксимации}

\noindent Исследование реализованных конфигураций на порядок аппроксимации состоит в
  отыскании предельной степени решения по времени ($k$), при котором исследуемый
  метод даёт точное решение.

\subsubsection{Описание задачи}

\begin{figure}[H]
  \centering
  \vspace{5mm}
  \includegraphics[scale=0.8]{pde/test/pbc/c1/img/geo*.png}
  \vspace{5mm}
  \caption{Засчётная область}
\end{figure}

\paragraph{Условия}
\indent
\vspace{1mm}

\noindent Решение - полином вида

\begin{align}
  u = x\sum_{p=0}^{k}(-1)^p t^p
\end{align}

\vspace{5mm}
\noindent Принимая значения коэффициентов 
  $\lambda = 1$, $\gamma = 0$, $\sigma = 2$, $\beta = 10$, получаем

\begin{align*}
  & f = \sigma x \sum_{p=1}^{k}(-1)^p p t^{p-1} \\
  & \theta|_{S_1} = -\sum_{p=0}^{k}(-1)^p t^p \\
  & \theta|_{S_3, S_5} = 0 \\
  & u_{\beta}|_{S_2} = x\sum_{p=0}^{k}(-1)^p t^p \\
  & u_{\beta}|_{S_4} = x\sum_{p=0}^{k}(-1)^p t^p + \frac{1}{\beta}\sum_{p=0}^{k}(-1)^p t^p \\
  & u_{\beta}|_{S_6} = x\sum_{p=0}^{k}(-1)^p t^p
\end{align*}

\subsubsection{Результат выполнения}

\begin{table}[H]
  \begin{tabularx}{\textwidth}{cYYY}
    \toprule
      $k$ & \multicolumn{3}{c}{$E(|u - u^*|)$}                              \\
    \cmidrule(lr){2-4}
          & Двухслойная схема   & Трёхслойная схема  & Четырёхслойная схема \\
    \midrule
      1   & 1.8910842e-11 & 1.0456255e-11 & 1.0443263e-11                   \\
      2   & 5.3892025e+00 & 2.0982637e-11 & 1.7016870e-11                   \\
      3   & 2.1751659e+02 & 1.3903611e+01 & 7.1946536e-10                   \\
      4   & 3.1009917e+03 & 9.5079913e+02 & 4.1711331e+01                   \\
    \bottomrule
  \end{tabularx}

  \captionsetup{justification=centering}
  \caption{Погрешность решения различных конфигурации в зависимости от степени полинома по времени}
\end{table}

% \begin{figure}[H]
%   \caption{График зависимости решения от времени в точке $(1, 0, 0)$}
% \end{figure}

\subsection{Тестирование на порядок сходимости}

\noindent Исследование реализованных конфигураций на порядок сходимости состоит в отыскании
  зависимости между приращением аргумента ($\Delta t$) и приращением погрешности ($E(|u - u^*|)$).

\subsubsection{Описание задачи}

\begin{figure}[H]
  \centering
  \vspace{5mm}
  \includegraphics[scale=0.8]{pde/test/pbc/c1/img/geo*.png}
  \vspace{5mm}
  \caption{Засчётная область}
\end{figure}

\paragraph{Условия}
\indent
\vspace{1mm}

\noindent Решение - периодическая функция вида

\begin{align}
  u = xsin(t)
\end{align}

\vspace{5mm}
\noindent Принимая значения коэффициентов 
  $\lambda = 1$, $\gamma = 0$, $\sigma = 2$, $\beta = 10$, получаем

\begin{align*}
  & f = 2xcos(t) \\
  & \theta|_{S_1} = -sin(t) \\
  & \theta|_{S_3, S_5} = 0 \\
  & u_{\beta}|_{S_2} = xsin(t) \\
  & u_{\beta}|_{S_4} = xsin(t) + 0.1sin(t) \\
  & u_{\beta}|_{S_6} = xsin(t)
\end{align*}

\subsubsection{Результат выполнения}

\begin{table}[H]
  \begin{tabularx}{\textwidth}{cYYY}
    \toprule
      $\Delta t$ & \multicolumn{3}{c}{$E(|u - u^*|)$}                             \\
    \cmidrule(lr){2-4}
                 & Двухслойная схема  & Трёхслойная схема  & Четырёхслойная схема \\
    \midrule
      128        & 1.9937299e+00 & 1.9653237e+00 & 1.8809388e+00                  \\
      64         & 4.0825225e+00 & 4.1217005e+00 & 4.1863085e+00                  \\
      32         & 2.9847026e+00 & 2.9443783e+00 & 2.9496760e+00                  \\
      16         & 2.6096649e+00 & 1.7627706e+00 & 8.5205136e-01                  \\
      8          & 2.9448070e+00 & 2.2266707e+00 & 1.2489770e+00                  \\
      4          & 6.3049236e+00 & 6.7250638e+00 & 7.1964954e+00                  \\
      2          & 3.3799714e+00 & 1.9801666e+00 & 8.9490448e-01                  \\
      1          & 1.7260509e+00 & 1.4720958e-01 & 8.8526785e-01                  \\
      0.5        & 8.6951389e-01 & 1.4709618e-01 & 1.5626740e-01                  \\
      0.25       & 4.3448994e-01 & 6.0461663e-02 & 1.7622318e-02                  \\
      0.125      & 2.2493731e-01 & 1.8818973e-02 & 2.0403619e-03                  \\
      0.0625     & 1.3608882e-01 & 3.6531704e-03 & 2.5589829e-04                  \\
      0.03125    & 2.6128831e-02 & 1.0842462e-03 & 1.1810640e-05                  \\
    \bottomrule
  \end{tabularx}

  \captionsetup{justification=centering}
  \caption{Погрешность решения различных конфигурации в зависимости от шага по времени}
\end{table}

\section{Проведенные исследования и выводы}

\section{Приложение}

\selectlanguage{english}
\inputminted[firstline=6, lastline=65]{c}{pde/src/fem/fem_lin_slv.c}
\inputminted[firstline=93, lastline=219]{c}{pde/src/fem/fem_lin_slv.c}
\inputminted[firstline=58, lastline=90]{c}{pde/src/fem/fem_lin_asm.c}
\inputminted[firstline=113, lastline=249]{c}{pde/src/fem/fem_lin_asm.c}

\end{document}
